\documentclass[../main.tex]{subfiles}
\begin{document}

\section{Results}
\label{sec:results}

\subsection{Overall Performance}
\label{sec:overall}

Figure~\ref{fig:overall-forest} summarises accuracy, Macro-F1, per-class F1, and Gwet's AC1
for all eight models across the 1{,}589 annotated cells, together with
study-clustered bootstrap 95\% confidence intervals (5{,}000 iterations,
resampling the 227 studies---the unit of clustering---with replacement to
preserve within-study correlation across the seven domains).
Overall accuracy ranged from 57.6\% (Gemini~2.5~Pro) to 62.5\%
(Claude~Haiku~4.5), a spread of less than five percentage points; the 95\%~CI
for each model spans approximately 5--6~pp, so all pairwise accuracy differences
fall within sampling uncertainty and should be treated as directional.
Macro-F1 confidence intervals similarly overlap across all eight models,
precluding reliable rank ordering based on this metric alone.
By contrast, Gemini~2.5~Flash and
DeepSeek~v4-Flash tied for the highest Macro-F1 (0.558) while
Claude~Haiku~4.5---the accuracy leader---ranked sixth (0.541).
The per-class breakdown reveals that the Macro-F1 gap is driven primarily by
High Risk and Unclear Risk detection: Claude~Haiku~4.5 achieved the highest
Low Risk F1 (0.743) but the lowest High Risk F1 (0.469) and Unclear Risk F1
(0.413) among all models.
Among the three metrics, Gwet's~AC1 shows the clearest separation at the
extremes: the 95\%~CI for Claude~Haiku~4.5 (AC1~$= 0.485$; CI~$[0.443, 0.526]$)
does not overlap with those of Gemini~2.5~Pro ($[0.344, 0.420]$) or
GPT-5.4 ($[0.349, 0.426]$), indicating that Haiku's higher chance-corrected
agreement is not attributable to sampling variation.
Consensus failures---cells where no model in the ensemble was correct---accounted
for 226 of 1{,}589 cells (14.2\%); the complementary any-correct rate of 85.8\%
confirms that the eight models make substantially non-redundant errors
(Section~\ref{sec:consensus-failures}).

\begin{figure}[ht]
  \centering
  \includegraphics[width=\linewidth]{all_reviews_overall_forest}
  \caption{Overall performance on RoB~1.0 assessment for the eight models
  ($n = 227$ studies; 1{,}589 cells).
  Each row is one model; columns show accuracy (\%), Macro-F1,
  F1 for each risk class (F1-LR, F1-HR, F1-UR), and Gwet's AC1.
  Filled circles are point estimates; horizontal bars are 95\% study-clustered
  bootstrap confidence intervals (5{,}000 iterations, resampling 227 studies
  with replacement).
  Dashed vertical lines mark the cross-model mean for each metric.
  Colours identify models consistently throughout the paper.
  Full numerical values are reported in Table~S\ref{tab:overall}.}
  \label{fig:overall-forest}
\end{figure}

\textbf{Metric divergence.}
The discrepancy between accuracy and Macro-F1 rankings reflects differential
handling of the class imbalance documented in Table~\ref{tab:class-dist}.
Figure~\ref{fig:overall-forest} makes the trade-off explicit:
models with higher Low Risk F1 tend to achieve higher accuracy but lower High Risk
and Unclear Risk F1, since Low Risk is the majority class in six of seven domains.
Macro-F1 penalises this behaviour by weighting each class equally, shifting
the ranking toward models with more balanced per-class detection.
Inspection of Figure~\ref{fig:overall-forest} also reveals that the rank orderings
for accuracy, AC1, and F1-LR are nearly identical across the eight models.
This co-linearity suggests that all three metrics are largely co-determined by
Low Risk performance: a model that reliably assigns Low Risk to majority-class
cells captures most of the variance in both accuracy and chance-corrected agreement,
since Low Risk prevalence (53.5\%) dominates both the observed agreement count and
the marginal distributions that enter the AC1 correction.
Macro-F1 disrupts this co-linearity by treating each class as equally important
regardless of prevalence.
However, equal class weighting does not reflect the empirical frequency of
judgements---High Risk accounts for only 17.4\% of cells
(Table~\ref{tab:class-dist})---so macro-F1 amplifies the contribution of the
rarest class beyond its share of the data.
A prevalence-weighted F1 would instead converge toward accuracy, whereas
macro-F1 should be interpreted as a measure of balanced class coverage:
how well a model performs across \emph{all three} nominal categories, independently
of how often each arises.

\subsection{Performance by RoB Domain}
\label{sec:domain}

% Figure: domain accuracy grouped bar chart
\begin{sidewaysfigure}
  \centering
  \includegraphics[width=\linewidth]{all_reviews_domain_grouped_bars_bootstrap}
  \caption{Per-domain accuracy by model (95\% bootstrap CIs; study-clustered,
  5{,}000 iterations, resampling 227 studies with replacement).
  Each group of eight bars corresponds to one RoB domain; colours identify models
  as shown in the legend.
  The dashed horizontal line marks 50\% (majority-class baseline).
  All analyses are exploratory.
  See Figure~\ref{fig:domain-metrics} for the breakdown by Macro-F1,
  per-class F1, and Gwet's AC1.}
  \label{fig:domain-bars}
\end{sidewaysfigure}

% Figure: domain metrics forest plot (5 panels)
\begin{figure}[ht]
  \centering
  \includegraphics[width=\linewidth]{all_reviews_domain_grouped_bars_metrics_bootstrap}
  \caption{Per-domain performance by model across five metrics: Macro-F1,
  F1 for Low Risk, Unclear Risk, and High Risk, and Gwet's AC1
  (95\% bootstrap CIs; study-clustered, 5{,}000 iterations,
  resampling 227 studies with replacement).
  Each panel shares the same domain grouping on the $x$-axis;
  diamonds are point estimates; vertical bars are 95\% CIs.
  The dashed horizontal line marks 0.5.
  Colours identify models consistently throughout the paper.}
  \label{fig:domain-metrics}
\end{figure}

Domain-level accuracy (Figure~\ref{fig:domain-bars};
Table~S\ref{tab:domain-full}) revealed two tiers of
difficulty.
\textbf{D1}~(sequence generation) and \textbf{D2}~(allocation concealment) were
the most accurately assessed domains, with mean accuracy of 75.3\% and 72.7\%,
respectively; all models exceeded 70\% accuracy on D1.
\textbf{D3}~(blinding of participants) showed high variance across models
(49.8\%--72.2\%), with Gemini~2.5~Flash and Claude~Haiku~4.5 outperforming
others by a margin of more than 5 percentage points.

\textbf{D7}~(other sources of bias) was the most challenging domain for all
models, with accuracy ranging from 24.2\% (Gemini~2.5~Pro) to 55.5\%
(Qwen~3.7-Plus)---a spread of 31 percentage points within a single domain.
Notably, models at opposite ends of the D7 accuracy range adopted diametrically
opposed prediction strategies: Gemini~2.5~Pro assigned High Risk to over 93\% of
D7 cells (recall 0.94 for High Risk; recall 0.18 for Low Risk), whereas
Qwen~3.7-Plus assigned Low Risk to 88\% of D7 cells (recall 0.88 for Low Risk;
recall 0.06 for High Risk).
Both strategies yielded similarly poor Macro-F1 in D7 ($\leq 0.37$ for both),
demonstrating that high accuracy on a single class does not translate to overall
competence.

Consensus failures were not uniformly distributed across domains: the
consensus failure rate was lowest for D2 (8.8\% of its 227 cells) and highest
for D4 (22.0\%), consistent with the greater class balance and annotation
difficulty in D4 relative to D2.

The domain-level pattern identifies two structurally distinct types of
difficulty.
\textbf{D4}~(blinding of outcome assessors) exhibits systematic, directionally
consistent errors: across all eight models, the large majority of errors are
lenient---models assign Low Risk when the ground truth is High Risk (see
Figure~\ref{fig:error-direction})---a pattern present across all intervention
types.
In drug trials in particular, models appear to conflate the objectivity of the
outcome measure with the adequacy of blinding of the outcome assessor,
defaulting to Low Risk when no explicit blinding failure is stated.
\textbf{D7}~(other sources of bias), by contrast, is characterised by
inter-model divergence: models adopt idiosyncratic strategies---some biased
toward High Risk, others toward Low Risk---so that errors are distributed roughly
equally between lenient and cautious directions.
The wide accuracy spread across models in D7 (31~pp) and its near-zero
chance-corrected agreement thus reflect not a shared systematic error but genuine
criterion ambiguity; the residual difficulty in D7 is unlikely to be resolved by
prompt refinement alone.

\subsection{Per-class Analysis}
\label{sec:perclass}

Figure~\ref{fig:domain-metrics} disaggregates Macro-F1 into its three constituent
classes.
Across all models, the ordering F1-LR~$\gg$~F1-UR~$\geq$~F1-HR holds as a
near-universal pattern, closely mirroring the prevalence ordering of the three
classes (53.5\%, 29.1\%, and 17.4\%, respectively; Table~\ref{tab:class-dist}).
Models detect each class roughly in proportion to how often they encounter it
in training and inference, producing a performance gradient that tracks
class frequency.
The sole exception is Claude~Haiku~4.5, which inverts the UR--HR ordering
(F1-HR~$= 0.469 >$ F1-UR~$= 0.413$), suggesting a different response
calibration: Haiku is more willing to assign High Risk when text signals a
methodological concern, at the cost of under-using Unclear Risk as an
epistemic hedge---a pattern consistent with its highest F1-LR (0.743) and
lowest F1-UR (0.413) in the entire benchmark.

\textbf{Low Risk} detection was strong and consistent across models in D1 and D2
(F1 $> 0.75$ for all models).
Performance deteriorated in D7, where most models reported Low Risk F1 below 0.60;
Gemini~2.5~Pro collapsed to F1~$= 0.29$ in D7 as a consequence of its High
Risk-biased prediction strategy described above.

\textbf{High Risk} detection was reliable only in D3 (F1: 0.47--0.79), the single
domain in which High Risk cases outnumbered Low Risk cases in the ground truth
(103 vs.\ 74; Table~\ref{tab:class-dist}).
The wide bootstrap confidence intervals for F1-HR visible in D1 and D2
in Figure~\ref{fig:domain-metrics} are a direct consequence of the small number
of High Risk cases in those domains (D1: $n = 5$; D2: $n = 10$ of 227 studies),
such that a single misclassified study shifts the F1 estimate substantially
under resampling; the point estimates themselves carry accordingly low precision
and should not be compared across models in those domains.
In D6~(selective outcome reporting), four models reported High Risk F1~$= 0.00$
(Gemini~2.5~Flash, GPT-5.4, Claude~Haiku~4.5, and Qwen~3.7-Plus), corresponding
to zero true-positive High Risk predictions across all 227 studies in that domain.
High Risk is the least prevalent class in D6 (25 of 227 cases, 11\%;
Table~\ref{tab:class-dist}), and these models predicted fewer than three High Risk
cases in D6 each, systematically defaulting to Low Risk or Unclear Risk in the
absence of explicit evidence of selective reporting.

\textbf{Unclear Risk} was the most consistently under-detected class.
Mean recall across all models fell below 0.20 in D3, D5, and D7.
The worst single result was Claude~Haiku~4.5 in D5: recall~$= 0.04$, meaning the
model correctly identified only 2 of the 45 Unclear Risk cases in D5.
The most accurately detected Unclear Risk domain was D2, where all models
achieved F1 $> 0.70$, consistent with the near-equal prevalence of Low Risk and
Unclear Risk judgements in that domain (113 vs.\ 104).

\textbf{Unclear Risk as abstention.}
The pattern of Unclear Risk under-detection reflects an asymmetric decision
behaviour.
In D3, D4, D5, and D7---domains where the ground-truth Unclear Risk rate ranges
from 20\% to 29\%---all models used Unclear Risk substantially less often than
the Cochrane reviewer, committing to definitive Low Risk or High Risk judgements
even in cases where the reviewer judged methodological information to be
insufficient.
In D2 and D6, the pattern reverses: models over-used Unclear Risk relative to
the ground truth, reflecting excessive hedging in the face of ambiguous reporting.
This bidirectional discrepancy indicates that models calibrate their
uncertainty threshold differently across domains and that reducing false
positives for Unclear Risk in the blinding domains (D3--D4) is a necessary
condition for improving High Risk recall there.

\subsection{Error Direction: Leniency and Caution}
\label{sec:error-direction}

% Figure: error direction heatmap
\begin{figure}[ht]
  \centering
  \includegraphics[width=\linewidth]{all_reviews_error_direction_leniency_optionC}
  \caption{Error direction by RoB domain and intervention type.
  Each cell shows the percentage of errors that are lenient
  (model assigns a lower-risk class than the ground truth)
  under an ordinal scale Low Risk $<$ Unclear Risk $<$ High Risk;
  $n$ = total errors in each cell.
  Values above 50\% indicate a lenient bias; values below 50\% indicate a
  cautious bias.}
  \label{fig:error-direction}
\end{figure}

Among cells where models erred, errors were not directionally neutral.
Figure~\ref{fig:error-direction} shows the percentage of errors that were
lenient across domains and intervention types.

\textbf{D4}~(blinding of outcome assessors) exhibited the most pronounced and
consistent leniency, with lenient errors exceeding 70\% in drug/medication
trials across almost all models.
This is interpretable: in pharmacological trials, the use of double-blind
placebo-controlled designs is common, and models may be defaulting to Low Risk
when such language is absent, rather than correctly inferring High Risk from
the lack of an explicit blinding protocol.

\textbf{D3}~(blinding of participants and personnel) showed a gradient by
intervention type: drug trials showed lower leniency (blinding is frequently
achieved), whereas device and procedural trials showed the highest leniency
rates, consistent with the fact that physical intervention concealment is
often impossible in those settings yet models continued to default to Low Risk.

\textbf{D1}~(sequence generation) and \textbf{D2}~(allocation concealment)
showed approximately balanced error directions across all models and intervention
types, consistent with their higher accuracy and the availability of explicit
methodological language in well-reported trials.
\textbf{D7}~(other sources of bias) showed no dominant direction, with models
diverging between lenient and cautious errors---confirming the heterogeneous
character of errors in this domain noted in Section~\ref{sec:domain}.

\clearpage
\subsection{Stratification by Intervention Type}
\label{sec:itype}

% Figure: domain accuracy by intervention type (4x1 panel)
\begin{figure}[H]
  \centering
  \includegraphics[width=\linewidth]{all_reviews_domain_grouped_bars_by_itype_bootstrap}
  \caption{Per-domain accuracy by model, stratified by intervention type
  (95\% bootstrap CIs; study-clustered, 5{,}000 iterations).
  Panels from top to bottom: drug/medication ($n = 104$),
  non-pharmacological intervention ($n = 75$),
  procedure ($n = 30$), and medical device ($n = 18$).
  Bootstrap CIs for the two smaller strata are wide and should be
  interpreted with caution; see Supporting Information~S4 for numeric estimates.}
  \label{fig:itype-bars}
\end{figure}

Mean Macro-F1 across all models ranked intervention types as follows:
non-pharmacological intervention (0.572, $n = 75$ studies),
drug/medication (0.538, $n = 104$),
medical device (0.531, $n = 18$), and
procedure (0.510, $n = 30$).
The 6.2-point gap between non-pharmacological interventions and procedures was
consistent across most models; no single model reversed this ordering.
Individual model rankings within each intervention type are shown in
Figure~\ref{fig:itype-bars} and provided in full in Table~S\ref{tab:itype-full}.

The interaction between intervention type and blinding assessment reveals a
specific failure mode.
In \textbf{D3}~(blinding of participants and personnel), accuracy for
drug/medication trials substantially exceeded that for device and procedural
trials; in the latter, the ground-truth High Risk prevalence is highest---because
physical concealment of the intervention is frequently impossible---yet models
assigned Low Risk to the majority of cells, producing leniency rates above 65\%
(Figure~\ref{fig:error-direction}).
In \textbf{D4}~(blinding of outcome assessors), the leniency bias reached its
maximum in drug trials (exceeding 80\%), where models appear to assume that
objective outcome measurement implies adequate assessor blinding.
These findings converge on a specific prompt deficiency: the current zero-shot
prompt does not distinguish between trial types when applying blinding criteria.
Models apply a ``blinding is feasible and likely'' default regardless of the
physical feasibility of concealment, an error that a single conditional rule
in the prompt (e.g., flagging as High Risk by default when concealment is
structurally impossible and no blinding mechanism is described) would
substantially reduce.

\subsection{Stratification by Clinical Domain}
\label{sec:clinical}

Among clinical domains with at least five studies, mean Macro-F1 across models
ranged from 0.425 (Infectious Disease, $n = 8$) to 0.680
(Orthopaedics \& Trauma, $n = 7$)---a 25.5-point spread
(Figure~\ref{fig:clinical-heatmap}).
The four highest-performing domains were
Orthopaedics \& Trauma (0.680),
Endocrine \& Metabolic (0.638, $n = 10$),
Rheumatology (0.637, $n = 5$), and
Wounds (0.628, $n = 5$).
The four lowest-performing domains were
Infectious Disease (0.425),
Reproductive \& Sexual Health (0.437, $n = 8$),
Genetic Disorders (0.473, $n = 9$), and
Mental Health (0.481, $n = 11$).
Complete results for all 24 Cochrane Review Groups are provided in
Table~S\ref{tab:clinical-full}.

These results are hypothesis-generating rather than confirmatory.
Group sizes of 5--11 studies at the extremes preclude reliable rank-ordering
across domains, and multi-label tagging makes groups non-exclusive (some reviews
are cross-listed under multiple Cochrane Review Groups).
The pattern is nonetheless qualitatively consistent with an interpretable
explanation: domains characterised by standardised reporting (orthopaedic
interventions and metabolic trials typically follow structured protocols with
explicit randomisation and blinding procedures) show higher Macro-F1, while
domains with heterogeneous or complex study designs (Infectious Disease,
Reproductive \& Sexual Health) present harder RoB characterisation
challenges for models trained on general text.
These observations motivate targeted future evaluation on larger domain-specific
corpora.

% Figure: clinical domain heatmap
\begin{sidewaysfigure}
  \centering
  \includegraphics[width=\linewidth]{all_reviews_clinical_domain_f1_heatmap}
  \caption{Macro-F1 by clinical domain and model (domains with $\geq 5$ studies;
  $n$ = number of studies in each domain, multi-label).
  Colour scale: 0--1.}
  \label{fig:clinical-heatmap}
\end{sidewaysfigure}

\subsection{Ground-Truth Uncertainty and Inter-Annotator Agreement}
\label{sec:iaa}

The benchmark treats published Cochrane annotations as the reference standard,
but those annotations carry their own measurement error: independent Cochrane
reviewer pairs achieve $\kappa = 0.24$--$0.37$ for non-D1 domains (Section~\ref{sec:conclusions}).
Figure~\ref{fig:agreement-matrix} places model performance in this context
by displaying the full $9 \times 9$ label-agreement matrix across Ground Truth
and all eight models.

\begin{figure}[H]
  \centering
  \includegraphics[width=\linewidth]{all_reviews_agreement_matrix}
  \caption{Label agreement matrix for Ground Truth and all eight models
  ($n = 1{,}589$ cells).
  Cell~$(i,j)$ = percentage of cells in which annotator~$i$ and
  annotator~$j$ assigned the \emph{same} label, irrespective of correctness.
  The diagonal is 100\% by definition.
  The white lines isolate the Ground Truth row and column.
  Colour scale: 55--100\% (Blues).}
  \label{fig:agreement-matrix}
\end{figure}

The Ground Truth row (first row) reproduces the known model accuracies
(57.6\%--62.5\%), representing agreement between each model and the
Cochrane reference.
All inter-model cells (the lower-right $8 \times 8$ block) exceed these values,
ranging from 63.1\% (Claude~Haiku~4.5 and Gemini~2.5~Pro) to 82.0\%
(Qwen~3.7-Plus and GLM-5).
These pairwise statistics are descriptive: no bootstrap confidence intervals
were computed for them.
Because cells from the same study share textual context and are correlated
across domains, the effective sampling unit is the study ($n = 227$), not the
cell ($n = 1{,}589$); each pairwise proportion therefore carries an approximate
95\% uncertainty of $\pm$6~pp ($\text{SE}\approx\sqrt{p(1-p)/227}$,
evaluated at $p\approx0.70$).
Differences smaller than this margin should not be over-interpreted; the
patterns noted below reflect the overall structure of the matrix rather than
fine-grained rankings between adjacent cells.
This systematic gap---models agree with each other more than they agree with
the ground truth---indicates that the eight models share systematic prediction
tendencies that diverge from the Cochrane standard in coordinated ways.
It does not imply that the models are collectively more reliable than the
ground truth; rather, it reflects shared training priors that cause
convergent errors.

The magnitude of the gap is interpretable in light of the known reference
noise: if Cochrane reviewer pairs themselves agree at $\approx$60--79\%
(the $\kappa = 0.24$--$0.37$ range, converted to expected agreement under
the observed class distribution), then the inter-model agreement range of
63--82\% is broadly comparable to or modestly above human inter-rater
reliability on this task.
The ground-truth column therefore represents one reference annotator among
many possible ones, and the model accuracy figures (57--63\%) should be
interpreted as lower bounds on the models' latent capability rather than
definitive measurements of their performance against a noise-free standard.

\subsection{Consensus Failures and Model Complementarity}
\label{sec:consensus-failures}

Of the 1{,}589 study--domain cells, 226 (14.2\%) were \emph{consensus failures}:
no model in the ensemble produced the correct judgement.
These cells may reflect genuine textual ambiguity, inconsistency in the Cochrane
annotation, or systematic limitations shared across all frontier models evaluated
here.
Consensus failures were concentrated in the most difficult domains: 50 (22.1\%)
occurred in D4 and 44 (19.5\%) in D5---the two domains with the most balanced
three-class distributions---corresponding to consensus failure rates of 22.0\%
and 19.4\% in those domains, respectively; D2 had the lowest consensus failure
rate (8.8\%).
Figure~\ref{fig:intersection} shows the pairwise intersection matrix:
cell~$(i,j)$ gives $P(A_i \cap A_j)$, the percentage of the 1{,}589 cells
in which \emph{both} model~$i$ and model~$j$ produced the correct judgement.
The diagonal entries report each model's overall accuracy ($P(A_i)$).

\begin{figure}[H]
  \centering
  \includegraphics[width=0.85\linewidth]{all_reviews_pairwise_intersection_heatmap}
  \caption{Pairwise intersection matrix.
  Cell~$(i,j)$ = $P(A_i \cap A_j)$: percentage of the 1{,}589 study--domain
  cells in which both models were correct.
  Diagonal entries (bold) = individual model accuracy.
  Colour scale: 40--65\% (YlOrRd); deeper red = greater overlap in correct cells.
  Upper triangle omitted by symmetry.}
  \label{fig:intersection}
\end{figure}

Off-diagonal intersection values range from 45.1\% (Claude~Haiku~4.5 and GPT-5.4)
to 53.5\% (Qwen~3.7-Plus and GLM-5), all substantially below the respective
diagonal entries (57.6\%--62.5\%).
The gap between diagonal and off-diagonal reflects that no two models
share all of their correct cells: a non-trivial fraction of each model's
successes are unique to that model.
Expressing overlap as the Jaccard similarity
$J(A,B) = P(A \cap B)\,/\,P(A \cup B)$---the fraction of cells correct
for \emph{at least one} model that are correct for \emph{both}---yields
values between 60.0\% and 78.1\% across the 28 pairs.
The highest Jaccard, Qwen~3.7-Plus and GLM-5 ($J = 78.1\%$, $P(A\cap B) = 53.5\%$),
indicates that nearly four in five cells that either model answers correctly
are answered correctly by both, confirming this as the most redundant pair.
The lowest Jaccard, Claude~Haiku~4.5 and GPT-5.4 ($J = 60.0\%$,
$P(A\cap B) = 45.1\%$), indicates that two in five jointly-correct cells
belong exclusively to one model---the highest non-redundancy in the ensemble
and a direct predictor of the high complementarity entry
visible in Figure~\ref{fig:complementarity}.

Figure~\ref{fig:complementarity} shows the full pairwise complementarity matrix:
cell~$(i,j)$ gives the percentage of the 1{,}589 cells in which model~$i$ erred
and model~$j$ was correct.
Values range from 7.4\% to 17.4\%, with all off-diagonal entries exceeding 7\%,
confirming that no two models in the ensemble are redundant.

\begin{figure}[H]
  \centering
  \includegraphics[width=0.85\linewidth]{all_reviews_pairwise_complementarity_heatmap}
  \caption{Pairwise complementarity matrix.
  Cell~$(i,j)$ = percentage of the 1{,}589 study--domain cells in which the
  row model erred and the column model was correct.
  Colour scale: 0--20\% (PuOr); purple = high complementarity,
  orange = low complementarity; grey diagonal = undefined.
  Full numeric values are reported in Table~S\ref{tab:pairwise}.}
  \label{fig:complementarity}
\end{figure}

The most salient feature of Figure~\ref{fig:complementarity} is the Claude~Haiku~4.5
column: it is the darkest column overall, meaning that when any other model errs,
Haiku recovers those cells at the highest rate (12.3\%--17.4\%).
This is the direct complement of its accuracy profile---Haiku leads on
overall accuracy and AC1 by concentrating predictions on the majority class,
which makes its correct cells less likely to be errors for other models,
but also makes its errors structurally distinct from theirs.
The two highest individual entries in the matrix---GPT-5.4~$\to$~Haiku (17.4\%)
and Gemini~2.5~Pro~$\to$~Haiku (17.3\%)---correspond to the two models with the
lowest overall accuracy; their larger error sets overlap least with Haiku's.
The same structural distinctiveness is visible in the agreement matrix
(Figure~\ref{fig:agreement-matrix}): Haiku has the lowest mean inter-model
agreement across the ensemble (67.9\%, versus 71.7\%--74.9\% for all other models),
driven by consistently low agreement with Gemini~2.5~Pro and GPT-5.4
(63--64\%)---the same dyads that top the complementarity column.
High complementarity and low overall agreement are two sides of the same coin:
because Haiku assigns labels differently from the ensemble cluster, it agrees
with other models less often irrespective of correctness, but those disagreements
are disproportionately cases where Haiku is right and the other model is wrong.
At the opposite extreme, the Qwen~3.7-Plus~$\leftrightarrow$~GLM-5 pair
shows the lowest bidirectional complementarity (7.4\% and 7.6\%), indicating
the most overlapping error profiles of any pair in the ensemble.
Grok~4.3 also shows low complementarity as a \emph{column} model for the
Qwen and GLM-5 rows (7.6\% and 8.9\%), suggesting that these three models
share a similar error structure that is not well covered by the others.
This pattern is reinforced by the agreement matrix (Figure~\ref{fig:agreement-matrix}):
the three highest pairwise inter-model agreements---Qwen~3.7-Plus / GLM-5 (82.0\%),
GLM-5 / Grok~4.3 (80.2\%), and Qwen~3.7-Plus / Grok~4.3 (77.6\%)---are all
concentrated within this trio; the trio's mean agreement (79.9\%) exceeds the
mean of cross-trio pairs by 7.2~pp.
Taken together, the matrix indicates that Claude~Haiku~4.5, Gemini~2.5~Pro,
and GPT-5.4 form the most complementary triad: their pairwise entries are
uniformly above 10\%, and combining any two of them recovers more than a
fifth of all cells that either individually misses.

\subsection{Inference Mode Comparison}
\label{sec:inference}

Table~\ref{tab:inference} compares per-domain and single-call inference modes
for the two models evaluated under both conditions.

\begin{table}[ht]
\centering
\caption{Per-domain vs.\ single-call inference mode comparison for
Gemini~2.5~Flash and DeepSeek~v4-Flash ($n = 227$ studies; 1{,}589 cells each).
F1 scores are computed globally across all cells (identical methodology to
Table~S\ref{tab:overall}); per-domain values reproduce the corresponding entries
in Table~S\ref{tab:overall}.
$\Delta$: single-call minus per-domain; 95\% paired bootstrap CI
(study-clustered, 5{,}000 iterations) shown in brackets.
$\dagger$: CI excludes zero (unadjusted for multiplicity; all analyses exploratory).}
\label{tab:inference}
\resizebox{\linewidth}{!}{%
\begin{tabular}{lcccccc}
\toprule
Model & Accuracy (\%) & Macro-F1 & F1-LR & F1-HR & F1-UR & Gwet's AC1 \\
\midrule
\multicolumn{7}{l}{\textit{Gemini~2.5~Flash}} \\
\quad Per-domain   & 60.7 & 0.558 & 0.718 & 0.474 & 0.484 & 0.437 \\
\quad Single-call  & 59.8 & 0.548 & 0.715 & 0.490 & 0.440 & 0.425 \\
\quad $\Delta$     & $-$0.9 & $-$0.010 & $-$0.003 & $+$0.016 & $-$0.044$^\dagger$ & $-$0.012 \\
\quad              & \footnotesize[$-$2.7, $+$1.0] & \footnotesize[$-$0.030, $+$0.011] & \footnotesize[$-$0.019, $+$0.014] & \footnotesize[$-$0.013, $+$0.051] & \footnotesize[$-$0.079, $-$0.009] & \footnotesize[$-$0.038, $+$0.015] \\
\midrule
\multicolumn{7}{l}{\textit{DeepSeek~v4-Flash}} \\
\quad Per-domain   & 60.3 & 0.558 & 0.704 & 0.480 & 0.489 & 0.433 \\
\quad Single-call  & 62.6 & 0.576 & 0.730 & 0.475 & 0.523 & 0.466 \\
\quad $\Delta$     & $+$2.3 & $+$0.018 & $+$0.026$^\dagger$ & $-$0.005 & $+$0.034 & $+$0.033 \\
\quad              & \footnotesize[$-$0.1, $+$4.7] & \footnotesize[$-$0.007, $+$0.044] & \footnotesize[$+$0.004, $+$0.047] & \footnotesize[$-$0.049, $+$0.041] & \footnotesize[$-$0.003, $+$0.072] & \footnotesize[$-$0.001, $+$0.066] \\
\bottomrule
\end{tabular}}
\end{table}

The overall picture from this two-model pilot comparison
(Table~\ref{tab:inference}) is one of small differences with wide uncertainty:
of 12~$\Delta$ values across both models and six metrics, ten have 95\%~CIs
that include zero.
The two whose CIs exclude zero are a drop in Unclear Risk F1 for
Gemini~2.5~Flash ($-0.044$; CI $[-0.079, -0.009]$, $\dagger$) and a gain
in Low Risk F1 for DeepSeek~v4-Flash ($+0.026$; $[+0.004, +0.047]$,
$\dagger$); both are class-specific and do not translate into differences
in accuracy, Macro-F1, or AC1, all of which remain within their respective
uncertainty intervals.
The accuracy differences ($-$0.9~pp for Flash, $+$2.3~pp for DeepSeek), while
pointing in opposite directions, are each compatible with zero change given the
bootstrap CIs ($[-2.7, +1.0]$ and $[-0.1, +4.7]$ respectively).

Domain-level results reinforce this picture.
D7~(Other Bias), the domain for which joint assessment was
hypothesised to confer the greatest benefit through shared contextual
reasoning, showed negligible differences under both modes ($\pm$1~pp)---
the most direct evidence that collapsing seven domain assessments into a single
call does not meaningfully improve the handling of the hardest domain.
The opposing D3 pattern (Flash: $-$0.4~pp; DeepSeek: $-$3.6~pp) and the
concentration of DeepSeek gains in D5 ($+$10.2~pp) and D2 ($+$4.4~pp)
suggest that whatever domain-level shifts occur are idiosyncratic to each
model rather than reflecting a systematic benefit of joint assessment.

Single-call outputs were systematically shorter (DeepSeek: 400$\to$214~chars;
Gemini~2.5~Flash: 530$\to$354~chars), with a corresponding decrease in hedging
expressions (e.g., ``unclear'', ``not reported'', ``cannot determine'') in
free-text justifications, consistent with a more compressed generation
regime under the single-call prompt.
The practical implications of this output compression appear limited and
model-dependent: it did not translate into a consistent performance benefit or
detriment across either model at the aggregate level.

\end{document}
